\section{Lemmas}

\begin{lemma}[subspace-imp-generator]
    \uses{def:LC_subspace,def:GeneratorMatrix}
    \label{lemma:subspace-imp-generator}
    Let $C$ be a linear code that is a $k$-dimensional subspace of $\mathbb{F}^{n}_{q}$.
    Then there exists a $k \times n$ matrix $G$ over $\mathbb{F}_q$ such that
    \[
        C = \{\mathbf{f}G \mid \mathbf{f} \in \mathbb{F}^{k}_{q}\}.
    \]
\end{lemma}

\begin{proof}
    It would follow first of all, that $C$ must have basis vectors since it is a vector space.
    Suppose we take $G$ to have rows that are just the basis vectors of $C$.
    It would further follow that by definition of dimension, there should be exactly $k$ such rows.
    From there we could argue that multiplication by all possible $f\in \mathbb{F}^{k}_{q}$ represents all possible linear combinations of the basis vectors and thus the vector space defined by them, i.e. $C$
\end{proof}

\begin{lemma}[generator-imp-subspace]
    \uses{def:LC_subspace, def:GeneratorMatrix}
    \label{lemma:generator-imp-subspace}
    Let $G \in \mathbb{F}^{k \times n}_{q}$ and let $C = \{\mathbf{f}G \mid \mathbf{f} \in \mathbb{F}^{k}_{q}\}$.
    Then $C$ is a subspace of $\mathbb{F}^{n}_{q}$.
\end{lemma}

\begin{proof}
    We can show that the zero vector $\vec{0}\in \mathbb{F}^{n}_{q}$ is in $C$ because $\vec{0}G=\vec{0}$
    We can show that $C$ is closed under addition by considering that for any arbitrary $a,b\in C$ it means that
        $\exists x,y\in \mathbb{F}^{k}_{q}$ such that $a=xG,b=yG$
        and from distributivity of matrix multiplication that $a+b=xG+yG=(x+y)G$ which should be in $C$ since
        $\mathbb{F}^{k}_{q}$ is closed under addition and thus $x+y\in \mathbb{F}^{k}_{q}$
    We can show that $C$ is closed under scalar multiplication by considering that for any arbitrary $a\in C$ that
        $\exists x\in \mathbb{F}^{k}_{q}$ such that $a=xG$ ad that therefore for any $d\in \mathbb{F}_{q}$
        $da=dxG$ and since $dx\in \mathbb{F}^{k}_{q}$ by its own closure under scalar multiplication it must follow $da\in C$
    These three conditions should be sufficient to prove that $C$ is a valid vector space.
    It is also *tRivIallY* true that because $C$ remains expressed in terms of $\mathbb{F}_{q}$ that $C\subseteq \mathbb{F}^{n}_{q}$
    Therefore $C$ is a subspace of $\mathbb{F}^{n}_{q}$
\end{proof}

\begin{lemma}[subspace-imp-parity-check]
    \uses{def:LC_subspace,def:ParityCheckMatrix,lemma:subspace-imp-generator}
    \label{lemma:subspace-imp-parity-check}
    Let $C$ be a linear code of dimension $k$ that is a subspace of $\mathbb{F}^{n}_{q}$.
    Then there exists a matrix $H \in \mathbb{F}^{(n-k) \times n}_{q}$ such that
    \[
        C = \{\mathbf{f} \in \mathbb{F}^{n}_{q} \mid H\mathbf{f}^{\intercal} = \mathbf{0}\}.
    \]
\end{lemma}

\begin{proof}
    Suppose we can use subspace-imp-generator to say that we know that $\exists G, \forall c \in C, \exists x\in \mathbb{F}^{k}_{q}, c=xG$
    We can extend this further to say that because the rows of $G$ are linearly independent since they are basis vectors and their quantity is $k$
    and they can be put into standard form $G=[I_k|P]$, (need to find if there's a good theorem to use)
    Then suppose we take $H=[-P^T|I_{n-k}]$ (I feel like handling these submatrices is gonna be hard ;-;)
    We can then see that $HG^T=I_k\cdot -P^T + I_{n-k}\cdot P^t= -P^T+P^t=0$
    From this we can get that since, by our defintion of $G$ that the aforementioned $x$ such that $c=xG$ can be used to show that
    $c\in C, \iff \exists x, c=xG \iff Hc^T=H(xG)^T=HG^Tx^t=0x^t=0$, meaning that $c\in C \iff Hc^T = 0$

\end{proof}

\begin{lemma}[parity-check-imp-subspace]
    \uses{def:LC_subspace, def:ParityCheckMatrix}
    \label{lemma:parity-check-imp-subspace}
    Let $H \in \mathbb{F}^{(n-k) \times n}_{q}$ and let $C = \{\mathbf{f} \in \mathbb{F}^{n}_{q} \mid H\mathbf{f}^{\intercal} = \mathbf{0}\}$.
    Then $C$ is a subspace of $\mathbb{F}^{n}_{q}$.
\end{lemma}

\begin{proof}
    We can show that the zero vector $\vec{0}\in \mathbb{F}^{n-k}_{q}$ is in $C$ because $H\vec{0}^T=\vec{0}$ trivially
    We can show that $C$ is closed under addition by considering that for any arbitrary $a,b\in C$ it means that
        $Ha^T=Hb^T=0$ and therefore by distributivity of matrix multiplication and transposition, $H(a+b)^T = Ha^T+Hb^T=0+0 = 0$ so $a+b\in C$
    We can show that $C$ is closed under scalar multiplication by considering that for any arbitrary $a\in C$ that
        $Ha^T=0$, then for arbitrary $d\in \mathbb{F}_{q}$, that $H(da)^T=dHa^T=0$ because we can commute scalars with matrices,
        giving us that $da\in C$
    These three conditions should be sufficient to prove that $C$ is a valid vector space.
    It is also true that since $C$ is defined in terms of elements of $\mathbb{F}^{n}_{q}$ that it is trivially a subset of $\mathbb{F}^{n}_{q}$
\end{proof}
